\section{Conclusiones}
El análisis experimental de ambos circuitos de corriente alterna permite extraer las siguientes conclusiones fundamentales:

\begin{enumerate}
    \item \textbf{Fenómeno de Resonancia (Parte A):} Se verificó experimentalmente la existencia de una frecuencia singular donde los efectos reactivos del capacitor y del inductor se cancelan mutuamente. En dicha condición, la impedancia se reduce al mínimo valor óhmico posible, maximizando la transferencia de corriente en el circuito. La selectividad del circuito quedó en evidencia al registrar caídas abruptas de tensión a ambos lados de la frecuencia central $f_0$.
    \item \textbf{Modelado de Componentes Reales (Parte B):} Los ensayos demostraron que los componentes comerciales se desvían de los modelos ideales. La bobina bajo estudio presentó una resistencia interna $r = 49.93\ \Omega$ comparable a la resistencia de carga ($R = 51\ \Omega$). Ignorar este factor geométrico hubiese inducido a un error severo en el cálculo de la inductancia. 
    \item \textbf{Validez de las Herramientas Fasoriales:} El uso de diagramas de fasores demostró ser una herramienta matemática precisa para describir la interacción de tensiones desfasadas en el tiempo, permitiendo resolver analíticamente parámetros no medibles de forma directa en el laboratorio.
\end{enumerate}

\end{document}